\section{Limitations}
\label{sec:limitations}

Our tasks are defect repairs of 15 to 20 minutes in an unfamiliar codebase.
Familiarity plausibly reduces resumption cost, since a developer who knows the
code rebuilds the suspended state faster, so our absolute numbers are likely an
upper bound on what a developer would experience in their own repository. We
have no reason to think familiarity changes the ordering of conditions, but we
did not test it.

The editor was purpose built. It looks and behaves like a mainstream tool, and
participants were given a training task, but it was not the tool any of them
use daily, and a small part of the measured lag is orientation in an unfamiliar
interface. This applies equally across conditions, so it inflates variance
rather than biasing the contrasts.

Builds were held at a fixed 42~s. Real builds vary, and a developer who cannot
predict the wait may disengage differently from one who can. Our follow-up at
6~s establishes that duration matters and locates a crossover, but the interior
of that range is unmapped.

The laboratory setting removes the social interruptions that dominate field
telemetry: colleagues, messages, and meetings. Those interruptions compete for
the same recovery capacity, and it is possible that in a working day saturated
with them the marginal cost of one more toast is smaller than we measure. It is
also possible that it is larger. Our design cannot distinguish these.

Finally, the sample is drawn from two meetups and one alumni list in a single
region and skews toward developers who volunteer for studies. The experience
distribution in Table~\ref{tab:demographics} is reasonable but the sample is
not representative of the profession, and the 12-participant follow-up in
Section~\ref{sec:discussion} is underpowered for anything beyond locating the
crossover approximately.
